\documentclass[11pt]{amsart}
\usepackage[T1]{fontenc}
\usepackage[utf8]{inputenc}
\usepackage{amssymb}
\usepackage[margin=1.05in]{geometry}
\usepackage[hidelinks]{hyperref}
\newcommand{\stamptime}{\the\numexpr\time/60\relax:\ifnum\numexpr\time-60*(\time/60)\relax<10 0\fi\the\numexpr\time-60*(\time/60)\relax}
\AddToHook{shipout/foreground}{\put(0.5\paperwidth,-\paperheight+0.55cm){\makebox[0pt]{\tiny\ttfamily unit/doc/\jobname.tex \ --- \ compiled \today, \stamptime}}}
\newtheorem{theorem}{Theorem}
\newtheorem{proposition}[theorem]{Proposition}
\newtheorem{lemma}[theorem]{Lemma}
\newtheorem{corollary}[theorem]{Corollary}
\theoremstyle{definition}
\newtheorem{definition}[theorem]{Definition}
\theoremstyle{remark}
\newtheorem{remark}[theorem]{Remark}
\DeclareMathOperator{\res}{res}
\newcommand{\sm}{\mathrm{span}}
\newcommand{\La}{\Lambda}
\title{Artin--Rees and the ACL certificate: an elementary account}
\date{28 October 2026}
\begin{document}
\maketitle

\noindent\textbf{Status.}
An elementary version of \texttt{artin-rees-acl.tex}: everything is
stated with polynomials and Laurent polynomials, without sheaves or
cohomology. All steps are checked by Unit (\texttt{unit/sag/artin-rees},
\texttt{unit/sag/acl-zmt}); file names are given in parentheses.
``Merely'' means: under propositional truncation (the object exists, but
is not part of the output).

\section{The two statements}

A ring $R$ is \emph{CFN} if (D) membership $v\in\sm(u_1,\dots,u_k)$ is
decidable, (C) the linear relations $\sum\lambda_iu_i=0$ between finitely
many elements are generated by finitely many of them, and (N) $R$ is
noetherian in the sense of Coquand--Persson. Discrete fields,
$\mathbb Z$, polynomial rings, quotients and localisations of CFN rings
are CFN; so is $A=k[b][a]$ for a discrete field $k$.

For $I=(a,b)$ and a submodule $P$, define $I^0P=P$ and: $x\in I^{m+1}P$
if $x=ax_1+bx_2$ with $x_1,x_2\in I^mP$.

\begin{theorem}[Artin--Rees]\label{thm:ar}
Let $R$ be CFN, $I=(a,b)$, $N=\sm(u_1,\dots,u_p)\subseteq R^r$. There
merely is $c$ such that, for every $m$: if $x\in I^{m+c}R^r$ and $x\in N$,
then $x=\sum_j\lambda_ju_j$ with every $\lambda_j\in I^m$.
\end{theorem}

In words: an element of $N$ that is very divisible by $I$ is a combination
of the generators of $N$ with coefficients that are almost as divisible,
the loss $c$ being independent of $m$
(\texttt{ArtinRees}, \texttt{PowSpan}).

\begin{theorem}[ACL]\label{thm:acl}
Let $A$ be a CFN gcd domain, $a,b\in A$, and
\[
 p=x+af+bg,\qquad q=y+au+bv\qquad(f,g,u,v\in A[x,y]).
\]
In every $A$-algebra $L$ containing $x,y$ with $p(x,y)=q(x,y)=0$ there
merely are $r,t\in L$ such that, for $s=1+ar+bt$, each of $s$, $sx$, $sy$
is a root of a monic polynomial with coefficients in $A$.
\end{theorem}

Modulo $I$, the equations say $x=y=0$: the solution is unique ``above
$a=b=0$''. The theorem turns this into an algebraic certificate. The
proof below constructs $s$ explicitly from $p,q$, except for finitely
many choices that are only merely given.

\section{The computational setting}

Homogenise: $P(X,Y,Z)=Z^dp(X/Z,Y/Z)$ and $Q=Z^eq(X/Z,Y/Z)$, so
\[
 P=XZ^{d-1}+aF+bG,\qquad Q=YZ^{e-1}+aU+bV .
\]

\textbf{Laurent polynomials with sign conditions.} Let $\La(n)$ be the
Laurent polynomials in $X,Y,Z$ over $A$, homogeneous of degree $n$
(exponents in $\mathbb Z$). For a set $E$ of variables, $\La_E(n)$ is the
part where only the variables in $E$ may have negative exponents. For
example $\La_Z(0)$ consists of the
$\sum c_{ij}X^iY^jZ^{-i-j}$ with $i,j\ge0$, i.e.\ of the polynomials in
$X/Z$ and $Y/Z$: it is $A[x,y]$ in disguise.

\textbf{Triples.} A \emph{$0$-triple} of degree $n$ is
$h=(h_X,h_Y,h_Z)$ with $h_X\in\La_X(n)$, $h_Y\in\La_Y(n)$,
$h_Z\in\La_Z(n)$. A \emph{$1$-triple} is
$(w_{XY},w_{XZ},w_{YZ})$ with $w_E\in\La_E(n)$. Put
\[
 \delta_0h=(h_Y-h_X,\ h_Z-h_X,\ h_Z-h_Y),\qquad
 \delta_1w=w_{XY}-w_{XZ}+w_{YZ}\in\La_{XYZ}(n),
\]
so that $\delta_1\delta_0=0$. Two triples are \emph{congruent},
$h\equiv h'$, when $h-h'=Pu+Qv$ componentwise, with $u,v$ triples of the
right shape and degrees $n-d$, $n-e$.

\textbf{Sections.} A $0$-triple $h$ is a \emph{section} if
$\delta_0h\equiv0$: its three components agree, modulo the equations, on
each overlap. Sections of degree~$0$ form a ring under componentwise
product, since
$(hh')_Z-(hh')_X=h_Z(h'_Z-h'_X)+h'_X(h_Z-h_X)$. The constant $(1,1,1)$
is a section (\texttt{CechSections}).

\textbf{Restriction.} For a $0$-triple $h$ of degree~$0$, $\res(h)\in L$
is obtained from $h_Z$ by setting $X/Z=x$, $Y/Z=y$. It is additive and
multiplicative, and kills $Pu+Qv$ because $p(x,y)=q(x,y)=0$. For $h$ of
degree $-1$, $\res(Zh)$, $\res(Xh)$, $\res(Yh)$ are defined, and
$\res(Xh)=x\,\res(Zh)$, $\res(Yh)=y\,\res(Zh)$, because $X=(X/Z)Z$.

So the plan is: \emph{find a section $h$ of degree $-1$ with
$\res(Zh)=1+ar+bt$}, then show that restrictions of sections are
integral over $A$.

\section{The proof in five steps}

\textbf{Step 1: an almost-section.} Let $z^*=(0,0,Z^{-1})$, of degree
$-1$. Then $\res(Zz^*)=1$, but $\delta_0z^*=(0,Z^{-1},Z^{-1})$, which is not
$\equiv0$: $z^*$ is not a section. It is, however, a section ``up to
arbitrarily high powers of $I$''. On the overlap $XZ$, where $X$ may be
inverted, put $\theta=X^{-1}Z^{1-d}$ and $Ef=-\theta(aF+bG)f$. Since
$P=XZ^{d-1}+aF+bG$, we get $f=Ef+\theta Pf$, hence
\[
 Z^{-1}=E^nZ^{-1}+P\,u_n,\qquad E^nZ^{-1}\in I^n\La_{XZ}(-1).
\]
This is the Laurent form of a simple fact: when $x$ is invertible,
$p=0$ gives $1=a(-f/x)+b(-g/x)\in I$. The same holds on $YZ$ with $Q$.
Hence $\delta_0z^*\equiv w_n$ with $w_n\in I^n$ for every $n$
(\texttt{FixedCochain}).

\textbf{Step 2: what stands between $w_n$ and $\delta_0$ of something
small.} If $w_n$ were $\delta_0c_n$ with $c_n\in I^n$, then $z^*-c_n$
would be a section with the same restriction as $z^*$ modulo $I^n$, and
we would be done. This is true only up to finitely many
\emph{obstructions}. Every Laurent polynomial $q$ of $\La_{XYZ}$
decomposes as $q=\pi(q)+\delta_1\sigma(q)$. Here $\pi(q)$ keeps the
monomials whose three exponents are all negative. Each other monomial
lies in some $\La_E$ with two variables, and $\sigma$ puts it there. If
$\delta_1w=0$, then $w=\delta_0\tau(w)$ for an explicit $\tau$
(\texttt{CechH1Finite}, \texttt{CechContraction}). Applying this to
$w_n$ gives
\[
 \delta_0z^*\equiv G\rho_n+\delta_0c_n,\qquad c_n\in I^n,\quad
 \rho_n\in A^r\text{ with entries in }I^{n-1},
\]
where $G\rho=\sum_t\rho_tm_t$ for fixed $1$-triples $m_1,\dots,m_r$, one
for each monomial with three negative exponents in a fixed degree
(\texttt{CechF1}). So the only obstruction is the vector
$\rho_n\in A^r$.

\textbf{Step 3: the obstructions that can be removed form a finitely
generated module.} Let
\[
 N=\{\rho\in A^r\mid G\rho\equiv\delta_0v\ \text{for some }v\}.
\]
By Step~2, $\rho_n\in N$, with $v=z^*-c_n$. Deciding whether
$G\rho\equiv\delta_0v$ has a solution amounts to linear equations over $A$
on finitely many coefficients. The only nonlinear input is the
description of the pairs $(\alpha,\beta)$ with $\alpha P+\beta Q=0$:
writing $P=D_0P'$, $Q=D_0Q'$ with $D_0=\gcd(P,Q)$, they are
$\alpha=-\gamma Q'$, $\beta=\gamma P'$ (the gcd hypothesis). This
expresses $N$ as the projection of the solutions of a finite matrix
$\mathcal M$ over $A$. By (C), $N$ is merely generated by
$u_1,\dots,u_p$, with explicit $b_j$ such that $Gu_j\equiv\delta_0b_j$
(\texttt{ACLRelations}).

\textbf{Step 4: the correction, by Artin--Rees.} Take $c$ from
Theorem~\ref{thm:ar} for $N=\sm(u_j)\subseteq A^r$, and $n=c+2$. Then
$\rho_n$ has entries in $I^{c+1}$ and lies in $N$, so
$\rho_n=\sum_j\lambda_ju_j$ with $\lambda_j\in I$. Put
\[
 h=z^*-c_n-\sum_j\lambda_jb_j .
\]
Then $\delta_0h\equiv G\rho_n-\sum_j\lambda_jGu_j=0$, so $h$ is a section.
Moreover $h-z^*\in I$, i.e.\ $h=z^*+am_1+bm_2$, and so
\[
 s=\res(Zh)=1+a\res(Zm_1)+b\res(Zm_2)=1+ar+bt,\qquad
 \res(Xh)=sx,\quad\res(Yh)=sy
\]
(\texttt{ACLCorrection}, \texttt{ACLCertificate}).

This is the step that needs Artin--Rees. Step~3 alone writes $\rho_n$ as
$\sum\mu_ju_j$ with arbitrary $\mu_j$. Then $h$ is still a section, but
$h-z^*$ is no longer in $I$, and $s$ need not be $\equiv1$ modulo $I$.

\textbf{Step 5: restrictions of sections are integral.} Every section $c$
of degree~$0$ is congruent to $\Phi(\kappa(c))=\kappa_0(c)\cdot1+\sum_i
\kappa_i(c)C_i$, for fixed sections $C_i$. Here
$\kappa(c)\in A^{s'}$ is computed from $c$ by the same $\pi,\sigma,\tau$
(\texttt{CechSections}). Apply the following lemma to
$f_n=\kappa(c^n)$ in $A^{s'}$.

\begin{lemma}[dependence]
If $A$ is CFN, every sequence $f_0,f_1,\dots$ in $A^{s'}$ merely has a
term $f_N\in\sm(f_0,\dots,f_{N-1})$ (\texttt{SeqDependence}).
\end{lemma}

This is noetherianity of $A^{s'}$ in its weakest form: the spans cannot
grow forever. Membership is decided by (D). From
$\kappa(c^N)=\sum_{i<N}x_i\kappa(c^i)$ and linearity of $\Phi$ we get
$c^N\equiv\sum_ix_ic^i$. Restricting, $\res(c)$ is a root of the monic
polynomial $T^N-\sum_ix_iT^i$. Applied to $c=Zh,Xh,Yh$, this proves
Theorem~\ref{thm:acl}.

\section{Remarks}

\begin{itemize}
\item Noetherianity of $A$ is used exactly twice: through Artin--Rees in
  Step~4, and through the dependence lemma in Step~5.
\item Every ``merely'' (the generators $u_j$, the bound $c$, the
  coefficients $\lambda_j$, the index $N$) is eliminated into the
  conclusion, which is itself a proposition. No choice principle and no
  excluded middle are used.
\item For readers who know the vocabulary: sections of degree $n$ modulo
  $\equiv$ are the global sections of $\mathcal O_W(n)$ on
  $W=\{P=Q=0\}$ in the projective plane over $A$. The obstructions of
  Step~2 are the classes in $H^1$, and Step~5 is the finiteness of
  $H^0$. None of this is used in the proof.
\end{itemize}
\end{document}
